\chapter{Ferritin (Blood)}
\section*{What Is This Test?}
Serum ferritin measures the concentration of ferritin, the primary intracellular iron storage protein, in the blood. Ferritin is a 450-kDa spherical protein composed of 24 subunits (heavy H-chain with ferroxidase activity, and light L-chain) forming a hollow shell that can store up to 4,500 iron atoms as a ferric oxyhydroxide-phosphate mineral core. Serum ferritin is in equilibrium with intracellular ferritin stores; 1 mcg/L of serum ferritin corresponds to approximately 8--10 mg of stored body iron. It is the single most useful laboratory test for assessing total body iron stores \cite{camaschella2015iron}. Ferritin is decreased in iron deficiency (the earliest and most specific indicator) and elevated in iron overload (hemochromatosis, transfusional siderosis). However, ferritin is also an acute phase reactant, synthesized by the liver and reticuloendothelial macrophages in response to inflammation, infection, malignancy, and liver disease. Ferritin is markedly elevated in hemophagocytic lymphohistiocytosis (HLH, >500 ng/mL, often >3,000--10,000+ ng/mL), Still disease (adult-onset Still disease, >1,000 ng/mL with marked hyperferritinemia), and COVID-19-associated hyperinflammatory syndrome. In these conditions, ferritin reflects macrophage activation and systemic inflammation, not iron overload.
\section*{Alternative Names}
\begin{itemize}
  \item Ferritin
  \item Serum Ferritin
  \item Iron Storage Protein
  \item SF
\end{itemize}
\section*{Principle}
Ferritin is measured by two-site sandwich chemiluminescent immunoassay (CLIA) or electrochemiluminescent immunoassay (ECLIA) using anti-ferritin monoclonal or polyclonal antibodies. The detection antibody is conjugated to a chemiluminescent label.
\section*{History}
Ferritin was first isolated from horse spleen in 1937 by Vaclav Laufberger, who gave the iron-rich protein its name from the Latin ``ferrum'' (iron). In the 1940s, Michaelis, Coryell, and Granick characterized its structure and iron core, establishing that ferritin is a protein shell surrounding a ferric oxide--phosphate mineral. The first immunoassays for serum ferritin were developed in the early 1970s \cite{addison1972immunoradiometric}, and the landmark 1975 paper by Jacobs and Worwood established that serum ferritin quantitatively reflects total body iron stores, making it the cornerstone of iron status assessment \cite{jacobs1975ferritin}. Radioimmunoassays were progressively replaced by automated chemiluminescent immunoassays in the 1990s and 2000s, which remain the standard platforms today. Recognition of ferritin as an acute phase reactant and of marked hyperferritinemia in hemophagocytic syndromes evolved over the 1990s--2010s, expanding its clinical interpretation beyond iron overload.
\section*{Why This Test Is Ordered}
\begin{itemize}
  \item Evaluation of microcytic anemia: low ferritin confirms iron deficiency as the cause
  \item Assessment of body iron stores in the workup of anemia of chronic disease, where iron deficiency may coexist with inflammation
  \item Diagnosis and monitoring of iron overload in hereditary hemochromatosis and transfusional siderosis, in which ferritin guides phlebotomy and chelation therapy \cite{bacon2011diagnosis}
  \item Monitoring of iron repletion during oral or intravenous iron therapy
  \item Evaluation of chronic kidney disease: ferritin and transferrin saturation guide the diagnosis of functional iron deficiency and erythropoiesis-stimulating agent therapy
  \item Investigation of markedly elevated ferritin (>1,000 ng/mL) for hemophagocytic lymphohistiocytosis, adult-onset Still disease, and severe COVID-19 \cite{cullis2018investigation}
  \item Routine assessment of iron status in pregnancy and in patients before surgery or bariatric procedures \cite{who2011serumferritin}
\end{itemize}
\section*{Primary vs Secondary Test}
Ferritin is a primary, first-line test for assessing total body iron stores and is the most sensitive and specific single marker of iron deficiency. It is always interpreted together with the CBC, serum iron, TIBC, and transferrin saturation; when inflammation confounds ferritin, the soluble transferrin receptor or bone marrow iron staining (the gold standard) are used.
\section*{How This Test Is Performed}
A blood sample is collected by standard venipuncture into a plain serum tube or a serum separator tube. The specimen is analyzed on an automated immunoassay platform using a two-site (sandwich) chemiluminescent or electrochemiluminescent immunoassay with anti-ferritin antibodies. If plasma is used rather than serum, the laboratory reference range may differ. Results are typically available within the same day.
\section*{What Preparation Is Needed}
Fasting is not required because ferritin is not affected by recent food intake. When ferritin is drawn as part of a full iron panel, an overnight fast is commonly recommended because serum iron shows significant diurnal variation. The sample should ideally be collected before iron supplementation begins so that baseline stores are measured.
\section*{Contraindications and Precautions}
\begin{itemize}
  \item There are no contraindications beyond those of routine venipuncture. Interpretive precautions:
  \item (1) Ferritin is an acute phase reactant, so concurrent inflammation, infection, malignancy, and liver disease elevate it and can mask true iron deficiency; in such patients a ferritin below 30 ng/mL, or below 100 ng/mL in chronic kidney disease, still strongly suggests depleted stores.
  \item (2) Recent blood transfusion and oral or intravenous iron therapy rapidly raise ferritin and obscure baseline stores.
  \item (3) An elevated ferritin alone does not distinguish iron overload from acute-phase inflammation; transferrin saturation and CRP are helpful adjuncts.
  \item (4) Severe liver injury releases intracellular ferritin, producing spuriously high values.
\end{itemize}
\section*{Risks}
Minimal risk. Slight pain or bruising at the venipuncture site. Rarely: hematoma, infection, or vasovagal syncope.
\section*{What Happens After the Test}
The ferritin result is interpreted alongside the CBC, serum iron, TIBC, transferrin saturation, and inflammatory markers. A low ferritin confirms iron deficiency and prompts iron replacement plus investigation of the underlying cause (gastrointestinal blood loss, malabsorption, or inadequate intake) \cite{guyatt1992laboratory}. An elevated ferritin with elevated transferrin saturation triggers evaluation for iron overload, including HFE gene testing and consideration of phlebotomy, whereas an elevated ferritin with normal transferrin saturation suggests an inflammatory or hepatic cause.
\section*{Understanding Results}
\subsection*{Decreased Ferritin (<15--30 ng/mL)}
Iron deficiency anemia (ferritin <15 ng/mL is essentially diagnostic of iron deficiency, 99\% specific). Pre-latent iron deficiency (depleted iron stores without anemia): ferritin 15--30 ng/mL, normal hemoglobin, and MCV. Ferritin <30 ng/mL in the presence of inflammation (elevated CRP) may be falsely elevated and mask iron deficiency. Ferritin <100 ng/mL in a patient with CKD/inflammation may still represent iron deficiency.
\subsection*{Elevated Ferritin}
Iron overload: Hereditary hemochromatosis (ferritin >300 ng/mL in men, >200 ng/mL in women), transfusional siderosis (ferritin >1,000+ ng/mL). Acute phase reactant: Infection, inflammation, malignancy (ferritin 500--5,000 ng/mL). Marked hyperferritinemia (>5,000--10,000+ ng/mL): HLH (macrophage activation syndrome, MAS), adult-onset Still disease, and severe COVID-19. Liver disease: Acute hepatitis, alcoholic liver disease, nonalcoholic fatty liver disease.
\section*{Normal Results}
Ferritin: Males 30--300 ng/mL, Females 20--200 ng/mL (premenopausal), 30--250 ng/mL (postmenopausal). Iron deficiency: <15 ng/mL (diagnostic). Iron deficiency with inflammation: may be normal or elevated (ferritin is an acute phase reactant); iron studies panel (serum iron, TIBC, transferrin saturation) and soluble transferrin receptor (sTfR) are needed. Newborns: markedly elevated ferritin at birth (25--200 ng/mL cord blood).
\section*{Follow-up Tests}
Serum iron, TIBC, transferrin saturation, soluble transferrin receptor (sTfR), CBC with RBC indices, HFE gene testing (hemochromatosis), liver function tests, and bone marrow biopsy with iron stain (gold standard).
\section*{References}
\bibliographystyle{unsrtnat}
\bibliography{references}
\clearpage
\section*{Regulatory Status and Authorization Date}
This chapter describes a test category, not a specific commercial product. FDA, CDSCO, EU IVDR, and MHRA approval or registration dates therefore cannot be assigned without the assay manufacturer, product name, instrument, and intended use. For laboratory-developed tests, an individual agency approval date may not exist. See the Preface for the interpretation of regulatory status.

\clearpage
